\documentclass[a4paper,11pt]{letter}

\usepackage[T1]{fontenc}
\usepackage[utf8]{inputenc}
\usepackage{lmodern}
\usepackage[hidelinks]{hyperref}
\usepackage{microtype}
\usepackage[margin=1in]{geometry}

\signature{Omar Iskandarani\\
Independent Researcher, Groningen, The Netherlands\\
ORCID: 0009-0006-1686-3961\\
Email: \href{mailto:info@omariskandarani.com}{info@omariskandarani.com}}

\address{Omar Iskandarani\\
Vinkenstraat 86A\\
9713 TK Groningen\\
The Netherlands}

\date{\today}

\begin{document}
    
    \begin{letter}{Editors\\ \textit{Foundations of Physics}}

\opening{Dear Editor,}

Please consider the enclosed manuscript,
\textit{Relational Time-of-Arrival as a Covariant Field Observable: From Event Currents to a Continuum Clock Limit},
for publication as a Research Article in \textit{Foundations of Physics}.

The manuscript proposes a covariant relational formulation of time-of-arrival in quantum field theory. Instead of introducing a universal time operator, it defines TOA as detector flux through a world-tube $\Sigma$ conditioned on a physical clock reading generated by an effective event-count current and its continuum coarse-grained limit.

The paper is intended as a contribution to the foundations of time observables, and is positioned explicitly in relation to operational TOA constructions and relational-clock frameworks. Unlike Page--Wootters-type constructions, the main object here is not a conditional state but a covariant detector functional defined from conserved currents and a clock-sector EFT. In that sense, the manuscript's contribution is structural: it formulates relational TOA directly at the field-observable level, with explicit UV-to-IR matching between discrete event counts and a continuum clock field, and with a worked $(1+1)$-dimensional example yielding a calculable arrival-time broadening formula.

I believe this is an appropriate submission for \textit{Foundations of Physics} because the manuscript is conceptually focused, technically explicit, and addressed to a recognizably foundational question: how arrival time should be defined when time is treated relationally rather than as an external parameter or universal operator.

The submission is original, has not been published previously, and is not under consideration elsewhere. A preprint version is available on Zenodo:
\href{https://doi.org/10.5281/zenodo.18050157}{10.5281/zenodo.18050157}.
There is no external funding and no conflict of interest.

Thank you for your consideration.

\closing{Sincerely,}
    
    \end{letter}

\end{document}